\scanpage{44}
b) Pour $a,\deleted{\unsure{b'}}a'\in A$, $b,b'\in B$, avec $a\ne a'$, $b\ne b'$, \deleted{\unsure{les ordres|les ord.}} $a,a'$ sont \deleted{\uline{contigus}} \inserted{adjacents} pour la structure \unsure{bi-circulaire|bi-circularisée} sur \inserted{l'ens à 4 éléments} $\{a,a',b,b'\}$ induite par celle de $S$, \deleted{i.e. ils} ($\Longleftrightarrow\{b,b'\}$ sont \deleted{contigus} \inserted{adjacents} $\Longleftrightarrow\{a,a'\}$ ne sont pas opposés $\Longleftrightarrow b,b'$ ne sont opposés

c) $B$ est contenu dans un segment composant de $S\setminus A$

c$'$) $A$ est contenu dans un segment composant de $S\setminus B$

% slanted
\margin{d) Quand $S$ est réalisé comme \inserted{partie} d'un disque \inserted{convexe} plongé dans un plan réel affine)\,: Soit $P_A$ la réunion des segments affines joignant deux à deux les points de $A$ (ou, \unsure{au choix|au besoin}, l'enveloppe convexe de $A$) alors $P_A\cap P_B=\varnothing$}

Quand il en est ainsi, parmi les couples $(\bar A,\bar B)$ \inserted{de segments} satisfaisant les conditions de a), il existe un plus petit\ldots

On dit que $A,B$ sont \uline{non panachés}

(\uline{Ex} Si \deleted{\unsure{$A$ et}} $A$ ou $B$ réduit à un pt. \inserted{(\unsure{plus généralt|plus généralmt} si $A$ ou $B$ est un segment}

Si $A$ et $B$ sont de cardinal 2, cel. signifie que, sur $A\cup B$ pour la bicircularisation induite, $A$ est un côté du carré $A\cup B$ (i.e. encore que $B$ est un côté). Dans le cas général, $A$ et $B$ sont non panachés ssi $\forall\, A'\in\unsure{P_2|\mathcal P_2}(A)$, $B'\in P_2(B)$, $A'$ et $B'$ sont non panachés.

% slanted
\margin{\deleted{e) \unsure{On a} $A,B\subset C$, \unsure{$C=$|un} segment \unsure{on a des} $\bar A,\bar B$ \unsure{ds|de} \unsure{st} $\bar A\cap\bar B=\varnothing$, \unsure{$=\ A$}}}

% slanted
\margin{\uline{Cor} De plus, \unsure{lorsque ces} conditions \unsure{sont} \unsure{sauf} dans le cas où $A\mathbin{\unsure{\sqcup|\cup}}B=S$ (\unsure{donc|alors} $A$ et $B$ des segments compl. l'un de l'autre), \unsure{quand il} existe \inserted{donc} un segment $C\supset A\cup B$, on peut prendre pour $\bar A,\bar B$ ds a) les \unsure{composantes de $A$ et $B$|composants de $A$ et $B$} $\subset C$.}
% archive sheet number: 43
